\section{Ocupación, exclusión y estadística finita}
\label{chap:ocupacion-estadistica}
\index{Pauli, principio de exclusión de}\index{CAR}
\index{Fermi--Dirac}\index{Bose--Einstein}

Las realizaciones nula y material proporcionan el soporte geométrico de la
ocupación. El sector fermiónico se construye sobre las relaciones canónicas de
anticonmutación (CAR), y el teorema de espín--estadística exige además la
representación operatoria declarada. Esta sección fija esas premisas antes de
derivar la idempotencia de la ocupación, las cardinalidades microcanónicas y
las distribuciones de Bose--Einstein y Fermi--Dirac con su límite clásico
común.

\subsection{Idempotencia de la ocupación en el álgebra de anticonmutación}

Sean \(a_i,a_i^\dagger\) operadores que satisfacen las relaciones canónicas
de anticonmutación (CAR)
\[
 \{a_i,a_j^\dagger\}=\delta_{ij},
 \qquad \{a_i,a_j\}=0,
\]
y sea \(n_i=a_i^\dagger a_i\).

\begin{theorem}[Principio de exclusión de Pauli]
\label{thm:principio-exclusion-pauli}
Se cumple \(n_i^2=n_i\), y por tanto
\(\Spec(n_i)\subseteq\{0,1\}\).
\end{theorem}
\begin{proof}
\[
 n_i^2=a_i^\dagger a_i a_i^\dagger a_i
 =a_i^\dagger(1-a_i^\dagger a_i)a_i
 =a_i^\dagger a_i,
\]
porque \(a_i^2=(a_i^\dagger)^2=0\). Todo autovalor de un idempotente
satisface \(\lambda^2=\lambda\) \cite{Pauli1925,Dirac1926}.
\end{proof}

La nilpotencia \((a_i^\dagger)^2=0\) impide la ocupación fermiónica doble del
mismo modo y conduce, al maximizar la entropía bajo las restricciones de
energía y número, a la ley de ocupación de Fermi--Dirac expuesta más adelante.
La completación simétrica y la ley de Bose--Einstein constituyen un estrato
distinto y se presentan por separado.

Las relaciones CAR son una hipótesis algebraica del teorema. Una monodromía
con cambio de signo bajo una rotación de \(2\pi\) es compatible con una
representación fermiónica, pero no constituye por sí sola una derivación de
CAR ni del teorema espín--estadística.
